\section{Worked Examples}
\label{sec:examples}

Three examples show how the framework changes what a study plans and reports.
The first two revisit our own recent work, because it is the work whose
strengths and omissions we know best; the third is a hypothetical study of a
kind that is becoming common. In each case we state the roles, the tier, and
the claim, and then what evidence the practices call for.

\subsection{A Speaker in a Robot-Led Wellbeing Intervention}
\label{sec:exA}

\paragraph{Study.} In a five-session intervention, a QTrobot whose dialogue
was generated by a GPT-3.5 model guided 21 university students through
cognitive reappraisal of emotional events \cite{laban2026robotled}. Loneliness
and perceived stress declined over sessions in mixed-effects models
\cite{laban2025what}.

\paragraph{Roles, tier, claim.} The model is a Speaker (\role{4}) and, because
it chooses which reappraisal move to make next, a Decider (\role{3}). Its
unreviewed output reaches people who are discussing distress, so the study is
Tier~III. The published claims are interaction-level and comparative across
time.

\paragraph{What the practices add.} The interaction-level evidence is what
HRI studies conventionally provide, and it is not weakened by anything that
follows. What it cannot show, by the argument of Section~\ref{sec:levels}, is
that the 560 responses the robot gave were faithful to the reappraisal
technique, accurate, and safe. A first component-level analysis was reported
separately: the robot's responses and the participants' disclosures were
compared with a corpus of 3{,}512 exchanges between help-seekers and licensed
therapists, and about 91\% of the disclosures made to the robot fell within
themes found in the therapy corpus, whereas only 28\% of therapy disclosures
fell within themes found in the robot corpus \cite{chiang2025talk}. That
asymmetry is itself informative about the narrower range of what a structured
robot-led session elicits. Reference-corpus alignment of this kind
(Section~\ref{sec:alignment}) establishes topical and semantic proximity to
professional practice at scale; it does not establish that any individual
response was appropriate, as the authors note.

Under the practices proposed here, the study would add: (i)~an audit in which
two raters with clinical training rate every robot response, or a
session-stratified random sample, on a rubric covering fidelity to the
intended technique, factual correctness, safety, honesty about the robot's
nature and capabilities, and persona consistency, with reliability reported
per dimension (\bpr{bp:human}, \bpr{bp:levels}, \bpr{bp:rubric}); (ii)~individual description of the
worst responses found (\bpr{bp:failures}); (iii)~if a model judge is used to extend the
audit to all responses, its validation against those raters on held-out
responses, relative to the raters' agreement with each other (\bpr{bp:validate},
\bpr{bp:judge}); (iv)~at system level, the latency distribution and the number of
times the supervising researcher intervened (\bpr{bp:levels}); and (v)~because GPT-3.5
snapshots have since been retired, an archived set of the dialogue contexts
so that a successor model can be compared with the original by replay
(Section~\ref{sec:replay}).

\subsection{A Representer and an Annotator in the Analysis of Self-Disclosure}
\label{sec:exB}

\paragraph{Study.} Disclosures that participants made to a human, a humanoid
robot, and a disembodied agent in three laboratory experiments were embedded
with a sentence-embedding model, clustered with $k$-means, and the clusters
were labelled and described by a language model; each description was then
checked by confirming that its embedding was closest to the centroid of its
own cluster \cite{chiang2025selfdisclosure}. The distribution of clusters did
not differ significantly by conversational partner, and the same pipeline was
used to relate disclosure themes to loneliness and stress in the intervention
described above \cite{laban2025what}.

\paragraph{Roles, tier, claim.} No model is in the interaction loop; the
robot and agent were controlled by a wizard. The embedding model is a
Representer (\role{2}) and the labelling model an Annotator (\role{7}), both
in the research pipeline. Their outputs define the dependent variable of the
inferential analysis, so the study is Tier~II. One of the claims, that themes
do not differ by partner, is a claim of absence.

\paragraph{What the practices add.} The pipeline already contains two
practices that we recommend generally. The automated description-to-centroid
check is an inexpensive test that a generated label is at least specific to
its cluster. More importantly, semantic similarity between partners was
computed with three embedding families, and the result differed sharply by
family: mean pairwise similarity ranged from about .27 to .49 with
sentence-transformer embeddings, and from about .79 to .94 with averaged
Word2Vec and BERT embeddings. Had only one family been used, the study would
have reported either dissimilarity or near-identity with equal confidence.
Triangulating across representations (\bpr{bp:embed}) exposed that the absolute
value of a cosine similarity is a property of the embedding space as much as
of the texts, which is why a fixed threshold such as .50 cannot carry the same
meaning across models.

The practices would add four things. First, an empirical null: the
distribution of similarity between randomly paired disclosures, within each
embedding space, against which observed similarity is expressed (\bpr{bp:embed}).
Second, a stability analysis of the clustering across seeds, bootstrap
resamples, and neighbouring values of $k$, since the elbow criterion selects
a single partition among several that are nearly as good (\bpr{bp:consistency}, \bpr{bp:embed}).
Third, human validation of the generated themes: an intrusion test in which
people identify the disclosure that does not belong to a cluster, and two
coders who assign a sample of disclosures to the model-generated themes, with
their agreement with each other and with the cluster assignment reported
(\bpr{bp:human}, \bpr{bp:coder}). Table~\ref{tab:kappa} suggests that with themes of unequal
size a sample of 200 or more is needed for the coefficient to be informative.
Fourth, for the claim that themes do not differ by partner, an equivalence
test or Bayes factor in place of a non-significant $\chi^2$
(Section~\ref{sec:equivalence}), since with 78 disclosures in the first
experiment the test had little power to detect a moderate difference.

\subsection{An Actor Handing Objects to People}
\label{sec:exC}

\paragraph{Study.} A hypothetical mobile manipulator in an assisted-living
facility is driven by a fine-tuned vision-language-action policy. Residents
ask for objects in their own words (``could you pass me my glasses, the ones
by the lamp''), and the robot fetches the object and hands it over.

\paragraph{Roles, tier, claim.} The policy is an Actor (\role{6}) that
subsumes perception and decision. It moves a physical arm near older adults,
so the study is Tier~III. The intended claim is of capability and safety.

\paragraph{What the practices ask for.} At component level, success is
estimated from a pre-specified number of physical trials per condition, with
exact intervals (Table~\ref{tab:cp}): thirty trials per condition are a
minimum for a capability claim, and more for a comparison between policies
(\bpr{bp:uncertainty}, \bpr{bp:trials}). Conditions are sampled along the axes that matter for the
interaction and that manipulation benchmarks do not vary: request phrasings
collected from residents instead of written by the researchers, seated and
standing recipients, occluded and outstretched hands, and bystanders in the
workspace (\bpr{bp:people}). Success is defined before data collection in terms of the
person (the object arrives in the hand in a safe orientation, without the
person having to lean or stand) and is adjudicated from video by two raters
whose agreement is reported; a vision-language model may be used as a success
detector only after it has been validated against those raters as a Judge
(\bpr{bp:human}, \bpr{bp:judge}). Failures are classified by type and by severity for the
person, so that a dropped spectacle case and an arm that advances towards a
face are not summed into the same failure rate (\bpr{bp:failures}).

At system level, the study reports the time from the end of the request to
the onset of motion, the behaviour of the safety layer that bounds speed and
separation independently of the policy, and every emergency stop and
teleoperated correction (\bpr{bp:levels}). Before residents are exposed, the policy is
stress-tested with ambiguous, contradictory, and adversarial requests
(Section~\ref{sec:redteam}). At interaction level, perceived safety, comfort,
and legibility are measured against a non-generative baseline, because task
success does not imply that being handed an object by this robot is tolerable
(\bpr{bp:people}). Finally, because the checkpoint will be updated, a small set of
physical regression scenarios and a larger simulated set are frozen so that
each update can be compared with its predecessor (Section~\ref{sec:replay}).
